\documentclass[11pt,a4paper]{article}
\usepackage[utf8]{inputenc}
\usepackage[T1]{fontenc}
\usepackage{amsfonts}
\usepackage[french]{babel}
\frenchbsetup{StandardLists=true} % à inclure si on utilise \usepackage[francais]{babel}
\usepackage{enumitem}
\usepackage{amssymb}
\usepackage{amsmath}
\usepackage[left=2cm,right=2cm,top=2cm,bottom=2cm]{geometry}
\usepackage{multirow}
\usepackage[dvipsnames]{xcolor}

\usepackage[T1]{fontenc}
\usepackage{fouriernc}
\usepackage{graphicx}

\usepackage{setspace}
\singlespacing

\usepackage{pstricks}
\usepackage{xkeyval}
\usepackage{pst-labo}


\usepackage{multicol}
\usepackage{caption}
\usepackage{fancyhdr}
\pagestyle{fancy}
\renewcommand\headrulewidth{1pt}
\fancyhead[L]{Lycée Denis Diderot }
\fancyhead[R]{Classe de TSN1}
\setcounter{page}{6}\fancyfoot[C]{page \thepage}
\fancyfoot[R]{\today}
\fancyfoot[L]{Alexis RUIZ}
\usepackage{multirow,tabularx,booktabs}
\newcommand{\cadreblanc}[1]{%
\medskip\framebox[\linewidth][c]{%
\rule{0mm}{#1mm}}\par
\medskip}

\usepackage{awesomebox}
\setlength{\aweboxvskip}{0mm}
\usepackage{pifont}
\usepackage{chemmacros}
\chemsetup{modules=all}
\setlength{\columnseprule}{.25mm}



\newcommand*\acocher{\quitvmode{\fboxsep0pt \fboxrule0.8pt \fbox{\vrule height1.8ex depth.1ex width0pt \vrule height0pt depth0pt width1.9ex }}}
\newcolumntype{R}[1]{>{\raggedleft\arraybackslash }b{#1}}
\newcolumntype{L}[1]{>{\raggedright\arraybackslash }b{#1}}
\newcolumntype{C}[1]{>{\centering\arraybackslash }b{#1}}

\newcommand{\defproblem}{\awesomebox[black]{2pt}{\rotatebox{90}{\large{\textbf{Problématique}}}}{black}}
\newcommand{\defconclusion}{\awesomebox[black]{2pt}{\rotatebox{90}{\Large{\textbf{Conclusion}}}}{black}}

\frenchbsetup{StandardLists=true}
\begin{document}



\begin{center}
\LARGE{COURS}
\end{center}




\hrule

\vspace{1cm}


\fbox{\textbf{Propagation d'un son}}\\

\cadreblanc{25}
\cadreblanc{25}
\vspace{2mm}

\fbox{\textbf{Caractéristiques d'une onde sonore}}
\begin{multicols}{2}
\textbf{Période}\\
\cadreblanc{35}
\textbf{Longueur d'onde}\\
\cadreblanc{34}
\end{multicols}

\fbox{\textbf{Pression acoustique et intensité acoustique}}
\begin{itemize}
\item Une onde sonore transporte de l'énergie, elle est caractérisée par l'intensité acoustique $I$ en W/$m^2$ qui représente la puissance moyenne transportée par l'onde par unité de surface. 
\item On mesure le niveau d'intensité acoustique $L$ en décibels (dB) avec un sonomètre, il est lié à l'intensité par la relation : \fbox{$L = 10 log \left(\dfrac{I}{I_0}\right)$} où $I_0$ est le seuil d'audibilité avec $I_0=10^{-12} W/m^2$.
\item \cadreblanc{14}

\end{itemize}

\fbox{\textbf{Caractéristique de l'oreille humaine}}\\

\cadreblanc{15}
\cadreblanc{12}
\end{document}